%%% Chapter 1 — Introduction
%%% UF-F10: this is the first of three chapters (intro / body / summary).
%%% UF-F11: heading hierarchy uses \chapter, \section, \subsection,
%%%         \subsubsection. The class auto-formats all five tiers.

\chapter{INTRODUCTION}\label{ch:intro}

Doctoral students at the University of Florida must submit their
dissertations through the Graduate Information Management System
(GIMS). Before submission, the dissertation must conform to a detailed
formatting specification maintained by the Graduate Editorial Office.
Common reasons for rejection include incorrect margins, omitted
required sections, non-functional hyperlinks, and improper use of
spacing or fonts.

This dissertation introduces a tool that helps students catch such
issues automatically before they trigger a rejection. The work is
motivated by three observations: (1) the official UF LaTeX template
already enforces most formatting rules by construction, but students
can and do override the template; (2) some properties --- for example,
the rendered margin on a page that includes a wide figure --- cannot be
determined from source alone; and (3) students near a submission
deadline benefit from a fast, offline tool that flags likely
problems with references back to the UF rule being violated.

If you are reading this as a soon-to-be doctor: the formatting
specification is finite, the rules are mechanical, and the validator
is on your side. You wrote the dissertation. The hard part is done.

\section{Background}

The UF Graduate School publishes its formatting specification at
\texttt{success.grad.ufl.edu/td/formatting/}. The UF IT Help Desk
publishes the LaTeX template at
\texttt{it.ufl.edu/helpdesk/graduate-resources/}. The current template,
updated for Fall 2025, requires LuaLaTeX with TeX Live 2025 to enable
PDF tagging for accessibility.

\section{Contributions}

This dissertation makes three contributions.

\subsection{A Two-Layer Validator Architecture}

A static analyzer that inspects both LaTeX source and rendered PDF.
Source-layer checks are precise and localizable to file and line.
PDF-layer checks confirm rendered-output truths the source layer
cannot predict.

\subsection{A Rule-by-Rule Catalog with Citations}

A documented mapping from UF formatting specification entries to
machine-checkable detection strategies, published as
\texttt{docs/uf-rules.md} so that the tool's checks can be audited
against the official specification.

\section{Organization of This Dissertation}

Chapter~\ref{ch:methods} describes the validator's architecture in
detail. Chapter~\ref{ch:conclusion} summarizes findings and discusses
future work. Appendix~\ref{app:rules} lists supplementary rule
identifiers, and Appendix~\ref{app:fixtures} describes the synthetic
broken inputs used to verify the validator's coverage.
